\chapter{Судалгааны сэдвийн онол, өнөөгийн түвшин}
\label{ch:literature}

Энэ бүлэгт эхлээд Монгол хэлний яриа таних аргуудыг тоймлоно (§\ref{sec:lit-asr}). Дараа нь сургуулийн баримтаас хариулт хайх аргуудыг авч үзнэ (§\ref{sec:lit-rag}).

%======================================================================
\section{Монгол хэлний яриа таних}
\label{sec:lit-asr}

\subsection{Яриа таних ба алдааны хэмжүүр}

Яриа автоматаар таних (\ac{ASR}, \ac{STT}) нь хүний ярианы аудиог текст болгох ажил юм. Чатботод хүн асуултаа ярихад эхлээд энэ алхам ажиллана. Энд гарсан алдаа дараагийн бүх алхамд дамждаг. Жишээ нь «хураамж» гэдэг үгийг буруу таньвал систем өөр мэдээлэл хайж эхэлнэ.

Яриа танилтын чанарыг ихэвчлэн үгийн алдааны хувиар (\ac{WER}) хэмждэг:
\begin{equation}
\mathrm{WER} = \frac{S + D + I}{N} \times 100\%
\label{eq:wer}
\end{equation}
Энд $S$ нь өөр үгээр сольсон үгийн тоо, $D$ нь орхисон үгийн тоо, $I$ нь илүү нэмсэн үгийн тоо, $N$ нь зөв бичвэрийн (reference) нийт үгийн тоо. Тэмдэгтийн алдааны хувь (\ac{CER}) ижил томьёогоор үсгийн түвшинд тооцогдоно. Нэмсэн үгийн тоо их байвал WER 100\%-иас давж болно.

WER дангаараа утгын алдааг бүрэн харуулдаггүй. Жишээ нь зөв бичвэр нь «хоёр мянга хорин тав», таньсан текст нь «2025» байлаа гэж үзье (зохиомол жишээ). Утга нь ижил. Гэвч үгээр тоолоход зөв бичвэрийн 4 үгээс нэгийг сольж, гурвыг орхисон тул WER $= 4/4 = 100\%$ болно. Эсрэгээр «2026» гэснийг «2020» гэж таньвал ганц үгийн алдаа. Гэвч асуултын утга бүрэн өөрчлөгдөнө. Иймээс нэр, тоо, оны алдааг тусад нь шалгах шаардлагатай (§\ref{subsec:lit-eval}).

\subsection{Whisper загвар}

Radford нар (2023) Whisper загварыг интернэтээс цуглуулсан 680 мянган цагийн аудио ба бичвэрийн хосоор сургасан \cite{radford2023whisper}. Үүний 117 мянган цаг нь англиас бусад 96 хэлний яриа юм. Хүн гараар шошголоогүй, интернэт дэх бичвэрийг шууд ашигласан тул энэ аргыг сул хяналттай сургалт (weak supervision) гэдэг. Загвар нь encoder--decoder бүтэцтэй Transformer бөгөөд нэг загвар яриа таних, хэл таних, англи руу орчуулах зэрэг хэд хэдэн ажил хийдэг. Гол давуу тал нь тухайн өгөгдөлд нэмэлт сургалт хийлгүйгээр (zero-shot) шууд хэрэглэж болдог явдал.

Гэвч сургалтын өгөгдөлд Монгол яриа маш цөөн байсан. Иймээс том хэмжээний Whisper large-v2 загварын Монгол хэлний WER нь FLEURS өгөгдөл дээр 110.5\%, Common Voice 9 өгөгдөл дээр 128.4\% гарсан \cite{radford2023whisper}. \shalgah{110.5 ба 128.4-ийг өгүүллийн хавсралтын хүснэгтээс тулгах} Өөрөөр хэлбэл алдааны тоо зөв бичвэрийн үгийн тооноос ч их байна. Зохиогчид урт аудиод үг давтах, орхих, огт хэлээгүй бичвэр зохиох алдаа гардгийг хязгаарлалт болгон дурдсан \cite{radford2023whisper}.

Whisper нь tiny-ээс large хүртэл хэд хэдэн хэмжээтэй. Энэ ажилд ашигласан small загвар 244 сая параметртэй \cite{openai_whisper_repo}.

\subsection{Whisper-ийг Монгол хэлэнд нэмэлт сургах}

Үндсэн Whisper загварын Монгол хэлний алдаа өндөр тул Монгол өгөгдлөөр нэмэлт сургасан (fine-tuned) загварууд гарсан. Тэдний нэг нь Hugging Face дээр нээлттэй байршуулсан \texttt{Cafet/whisper-small-mongolian} загвар юм \cite{cafet_whisper_small_mn}. Энэ нь Whisper small-ийг Монгол хэлээр нэмэлт сургасан хувилбар. Загварын тайлбар хуудсанд сургалтын тохиргоо, шалгалтын WER бичигдсэн. Гэхдээ энэ нь рецензтэй судалгааны өгүүлэл биш. Мөн тэнд бичсэн WER манай сургуулийн асуултад ижил байна гэсэн баталгаа биш.

Mengke нар (2025) Халх Монгол хэлний өгөгдөл цөөн байх асуудлыг шийдэх аргыг санал болгосон \cite{mengke2025khalkha}. Common Voice 17-д Халх Монгол хэлний сургалтын өгөгдөл ердөө 3 цаг байдаг. Иймээс зохиогчид Өвөр Монголын Чахар аялгууны 147 цагийн өгөгдлийн бичвэрийг кирилл үсэг рүү хөрвүүлсэн (transliteration). Эхлээд энэ өгөгдлөөр, дараа нь Халхын өгөгдлөөр хоёр шаттай нэмэлт сургалт хийсэн. Whisper large-v3 загварын WER нь шууд хэрэглэхэд 92.60\%, зөвхөн Халхын өгөгдлөөр нэмэлт сургахад 43.02\%, хоёр шаттай аргаар 31.45\% болж буурсан. Хоёр өгөгдлийг шууд нийлүүлж сургахад 49.95\% гарсан. Сургалтад 48 GB санах ойтой A6000 GPU ашигласан \cite{mengke2025khalkha}.\footnote{Зохиогчид §3.2-т Whisper-ийг CTC алдагдлын функцээр сургасан гэж бичсэн. Энэ нь Whisper-ийн ердийн encoder--decoder сургалтаас өөр тул энэ хэсгийг нягтлах шаардлагатай. \shalgah{өгүүллийн §3.2-ыг дахин уншиж тулгах}}

Энэ судалгаанаас гурван сургамж гарна. Нэгдүгээрт, Whisper-ийг Монгол хэлэнд мэдэгдэхүйц сайжруулах боломжтой. Хоёрдугаарт, үүнд олон цагийн аудио--бичвэрийн хос, хүчтэй GPU, тусдаа шалгах өгөгдөл шаардлагатай. Гуравдугаарт, тэдний шалгалт уншсан яриа (Common Voice) дээр хийгдсэн тул сургуулийн чөлөөт асуултад ижил дүн гарна гэж хэлэх боломжгүй.

\subsection{Gemini 3.5 Transcribe}

Gemini 3.5 Transcribe нь Google-ийн Gemini \ac{API}-аар ажилладаг, аудио файлыг текст болгох зориулалттай үүлэн үйлчилгээ юм \cite{google_transcribe_model}. Албан баримтад Монгол хэл (\texttt{mn-MN}) дэмжих хэлний жагсаалтад орсон. Хэрэглэгч хэлийг урьдчилан зааж (language hint), үгчилсэн (verbatim) горимоор бичүүлж болно \cite{google_transcribe_api}.

Энэ үйлчилгээг ашиглахад \ac{API} түлхүүр, интернэт холболт шаардлагатай. Аудио Google-ийн сервер рүү илгээгддэг. Үнэгүй түвшинд (free tier) илгээсэн өгөгдлийг Google бүтээгдэхүүнээ сайжруулахад ашиглаж, хүн хянаж болно гэж нөхцөлд заасан \cite{google_terms}. Иймээс хувийн мэдээлэлтэй аудиог илгээхэд болгоомжлох хэрэгтэй. Төлбөртэй түвшний ойролцоо үнэ минутад 0.005 ам.доллар орчим байна (2026.10.29-ний байдлаар) \cite{google_pricing}.

Google-ийн өөр нэг загвар \texttt{gemini-3.5-live-translate-preview} нь яриаг өөр хэл рүү шууд орчуулах зориулалттай \cite{google_live_translate}. Түүний оролтын Монгол текстийг авч болох ч үндсэн зорилго нь орчуулга. Иймээс энэ ажилд түүнийг Transcribe-аас тусад нь авч үзсэн (§\ref{sec:res-demo}).

Gemini 3.5 Transcribe-ийн Монгол хэлний алдааг хэмжсэн рецензтэй судалгаа энэ ажлын хүрээнд олдоогүй. Баримтад хэл «дэмждэг» гэж бичсэн нь манай ярианд сайн таних баталгаа биш. Иймээс өөрсдийн бичлэг дээр хэмжих шаардлагатай болсон.

\subsection{Хоёр аргын харьцуулсан тойм}

\tref{tab:asr_compare}-т хоёр аргыг онолын талаас нь харьцуулав. Хамгийн гол ялгаа нь Whisper-ийг сайжруулахын тулд нэмэлт сургалт хийх шаардлагатай байхад Gemini-г API түлхүүрээр шууд ашиглаж болох явдал юм. Харин Gemini-д аудио гадагш гардаг, загварыг өөрөө сургах боломжгүй.

\begin{table}[H]
\centering
\caption{Whisper ба Gemini 3.5 Transcribe-ийн шинж чанарын харьцуулалт}
\label{tab:asr_compare}
\small
\begin{tabular}{|p{3.4cm}|p{5.1cm}|p{5.1cm}|}
\hline
\textbf{Шинж} & \textbf{Whisper (Cafet small)} & \textbf{Gemini 3.5 Transcribe} \\
\hline
Хаана ажиллах & Өөрийн компьютер (CPU эсвэл GPU) & Google-ийн үүлэн үйлчилгээ \\
\hline
Интернэт & Шаардлагагүй & Шаардлагатай \\
\hline
Аудио хаашаа очих & Компьютер дотроо үлдэнэ & Google-ийн сервер рүү очно \\
\hline
Монгол хэлэнд сайжруулах & Нэмэлт сургалт: олон цагийн өгөгдөл, GPU \cite{mengke2025khalkha} & Загварыг сургах боломжгүй; хэл заах, горим сонгох \\
\hline
Ашиглахад хэрэгтэй & Загварын файл, тооцоолох нөөц & API түлхүүр \\
\hline
Зардал & Тоног төхөөрөмж, цахилгаан & Үнэгүй квоттой; төлбөртэй нь минутад $\approx$0.005 ам.доллар \\
\hline
Нийтлэгдсэн Монгол хэлний хэмжилт & large-v2: FLEURS дээр WER 110.5\% \cite{radford2023whisper}; нэмэлт сургасан large-v3: 31.45\% \cite{mengke2025khalkha} & Рецензтэй хэмжилт олдоогүй \\
\hline
\end{tabular}
\end{table}

\subsection{Яриа танилтыг үнэлэхэд анхаарах зүйл}
\label{subsec:lit-eval}

Өмнөх судалгаанаас яриа танилтыг үнэлэхэд дөрвөн анхаарах зүйл гарсан.

\begin{enumerate}
    \item \textbf{Нэг арга бүх хэлэнд хамгийн сайн биш.} Jimerson нар 11 хэл дээр хэд хэдэн архитектурыг харьцуулахад нэг ч арга бүх хэлэнд тэргүүлээгүй \cite{jimerson2023unhelpful}\unshih. Иймээс сонголтоо өөрийн өгөгдөл дээр хийх хэрэгтэй.
    \item \textbf{Текстийг хэвийн болгох (normalization) нь дүнг өөрчилнө.} Цэг тэмдэг, тэмдэгтийг арилгах дүрэм зарим хэлэнд WER-ийг хиймлээр бууруулдаг \cite{manohar2024normalization}\unshih. Иймээс хоёр аргыг яг ижил дүрмээр тоолох ёстой.
    \item \textbf{Нэр, тооны алдаа дараагийн ажлыг эвддэг.} Яриа танилтын алдаа нэр томьёог таних ажилд хүчтэй нөлөөлдөг \cite{szymanski2023ner}\unshih. Манай чатботод нэр, тоо, он хайлтын гол түлхүүр болно.
    \item \textbf{Яриагүй хэсэгт бичвэр зохиох.} Whisper заримдаа огт хэлээгүй өгүүлбэр зохиодог \cite{koenecke2024careless}\unshih. Иймээс чимээгүй бичлэгээр шалгах хэрэгтэй.
\end{enumerate}

%======================================================================
\section{Мэдээлэл хайх ба RAG}
\label{sec:lit-rag}

\subsection{RAG-ийн үндсэн санаа}

Lewis нар (2020) \ac{RAG} аргыг санал болгосон \cite{lewis2020rag}\unshih. Энэ арга хоёр хэсгээс бүрдэнэ. Эхний хэсэг нь хайгч (retriever). Хайгч асуултад тохирох баримтын хэсгийг том сангаас олно. Зохиогчид Wikipedia-г 100 үгтэй 21 сая орчим хэсэг болгон хуваасан. Хоёр дахь хэсэг нь хариулт бичигч (generator). Тэр асуулт болон олсон хэсгийг уншиж хариулт бичнэ.

Энэ аргын гол санаа нь мэдлэгийг загварын параметр дотор биш, гадаад баримтын санд хадгалах явдал юм. Иймээс хоёр давуу тал гарна. Нэгдүгээрт, хариулт аль баримтаас гарсныг харуулах боломжтой. Хоёрдугаарт, загварыг дахин сургахгүйгээр баримтын санг шинэчилж мэдлэгээ өөрчилж болно. Сургуулийн мэдээлэл жил бүр өөрчлөгддөг тул энэ шинж манай ажилд чухал.

Жишээ авъя. Хэрэглэгч «Бүртгэлийн хураамж хэд вэ?» гэж асуув. Хайгч сургуулийн баримтаас бүртгэлийн хураамжийн тухай хэсгийг олно. Дараа нь хариулт тэр хэсэгт тулгуурлан гарна. Хэрэглэгч хариултын хажууд эх баримтыг нь харна.

Манай системд хайх хэсэг (R) нь гуравдугаар алхам буюу BM25 ба E5 хайлт юм (§\ref{sec:architecture}). Хариулт гаргах хэсэг (AG) нь дөрөвдүгээр алхам. Гэхдээ манай систем шинэ өгүүлбэр бичихгүй. Олсон хэсгүүдээс тохирох 1--3 хэсгийг дүрмээр сонгоод, эхийн үгээр нь харуулдаг. Давуу тал нь эх баримтад байхгүй зүйл зохиогдохгүй. Сул тал нь хэд хэдэн хэсгийн мэдээллийг нэгтгэж товч хариулж чаддаггүй. Хариулт сонгоход LLM ашиглах эсэхийг W08-д ижил асуултаар харьцуулна.

\bichne{Өгүүллийг уншсаны дараа энэ хэсгийг өөрийн үгээр засаж, RAG-Sequence ба RAG-Token хувилбарын ялгааг нэмнэ.}

\subsection{Үгээр хайх: BM25}
\bichne{W04 --- \cite{robertson2009bm25}\unshih: ховор үг яагаад их оноо авдаг, давтагдсан үгийн оноо яагаад хязгаартай вэ.}

\subsection{Утгаар хайх: вектор ба E5}
\bichne{W07 --- \cite{wang2024me5}\unshih: текстийг вектор болгох, косинусын ойролцоо, Монгол хэлний хязгаар.}

\subsection{Хоёр хайлтыг нэгтгэх: RRF}
\bichne{W05 --- \cite{cormack2009rrf}\unshih: $\sum 1/(k+\text{байр})$ томьёо; өгүүлэлд $k=60$, манайд $k=30$.}

\subsection{Мэдлэгийн мод ба бүтэцтэй хайлт}

Энгийн RAG бүх баримтын хэсгийг нэг жагсаалт мэт хайдаг. Сүүлийн үеийн аргууд мэдлэгээ мод эсвэл граф хэлбэрээр зохион байгуулдаг. RAPTOR нь ойролцоо хэсгүүдийг бүлэглэж, LLM-ээр хураангуйлан доороос дээш мод үүсгэдэг \cite{sarthi2024raptor}\unshih. GraphRAG нь текстээс нэр, холбоог LLM-ээр гаргаж граф үүсгэдэг \cite{edge2024graphrag}\unshih. Эдгээр нь маш том текстэд зориулагдсан бөгөөд мод, графаа LLM-ээр бүтээдэг.

Манай мэдлэгийн мод хүний гараар хийсэн JSON файл юм. Баримтын сан 200 орчим хэсэгтэй жижиг. Иймээс модыг хураангуй бичихэд бус, дараах дөрвөн зорилгоор ашиглана (§\ref{sec:knowledge-tree}):
\begin{enumerate}[nosep]
    \item асуултыг тохирох сэдэв рүү чиглүүлэх;
    \item тухайн сэдвийн хэсгүүдийг хайлтад дээш эрэмбэлэх;
    \item багш, өрөө, анги, хуваарийн хүснэгтийг сэдэвтэй холбож шууд хариулах;
    \item мэдээлэл байхгүй сэдвийг таньж, «мэдээлэл алга» гэж хэлэх.
\end{enumerate}
\bichne{W09-д RAPTOR, GraphRAG-ийг уншсаны дараа энэ харьцуулалтыг гүнзгийрүүлнэ.}

\subsection{Хариугүй асуулт}
\bichne{W03 --- \cite{rajpurkar2018squad2}\unshih: сэдэв ойр боловч хариугүй асуулт яагаад хэцүү, систем хэзээ «олдсонгүй» гэх ёстой вэ.}

%======================================================================
\section{Төсөөтэй системүүд}
\bichne{Их сургуулийн элсэлтийн чатботод дүрэм ба RAG-ийг хослуулсан URAG \cite{nguyen2025urag}\unshih зэрэг ажлыг уншиж, манай системтэй хүснэгтээр харьцуулна.}
